\paragraph{斜入射要连续取两次法向投影}
令 $I$ 是按垂直于光束的面积定义的入射光强，$\theta$ 从表面法线量起。实际表面单位面积每秒接收的能量是 $I\cos\theta$；每单位能量携带的法向动量为 $\cos\theta/c$。因此完全吸收表面的法向辐射压为
\[
 p_{\perp,\rm abs}=\frac I c\cos^2\theta,
\]
完全镜面反射时法向动量反号，故 $p_{\perp,\rm refl}=2I\cos^2\theta/c$。吸收表面还可能承受切向动量通量，不能把法向压强误作总作用力大小。
如果题目给的是按实际表面面积计的辐照度 $\mathcal E=I\cos\theta$，同一结果应写为 $p_{\perp,\rm abs}=\mathcal E\cos\theta/c$，否则会重复投影。这里假定真空中的单向光束；各向同性辐射场的角平均给出 $p=u/3$，不能与单向束的 $p=u$ 混用。
